\documentclass[11pt,a4paper]{article}

\usepackage[utf8]{inputenc}
\usepackage[T1]{fontenc}
\usepackage{lmodern}
\usepackage{microtype}
\usepackage{amsmath, amssymb}
\usepackage{mathtools}
\usepackage{booktabs}
\usepackage{array}
\usepackage{tabularx}
\usepackage{multirow}
\usepackage{xcolor}
\usepackage{hyperref}
\usepackage{geometry}
\usepackage{listings}
\usepackage{enumitem}
\usepackage{titlesec}
\usepackage{abstract}

\geometry{
  top=2.5cm, bottom=2.5cm,
  left=2.5cm, right=2.5cm
}

\hypersetup{
  colorlinks=true,
  linkcolor=blue!60!black,
  urlcolor=blue!70!black,
  pdftitle={PLAA: Persistent Latent Adaptive Architecture},
  pdfauthor={Xiao Ouyang}
}

\lstset{
  basicstyle=\small\ttfamily,
  backgroundcolor=\color{gray!10},
  frame=single,
  framesep=4pt,
  breaklines=true,
  columns=fullflexible,
  keepspaces=true,
  xleftmargin=0.5em,
  xrightmargin=0.5em
}

\newcommand{\R}{\mathbb{R}}
\newcommand{\St}{S_t}
\newcommand{\Stp}{S_{t+1}}
\newcommand{\hs}{h_s}

\title{
  \textbf{PLAA: Persistent Latent Adaptive Architecture}[0.4em] \large Structural Symmetry Breaking Without Functional Coupling:\\ \large Latent Modulation as Training-Time Scaffold in Frozen LLMs }  \author{ Xiao Ouyang\thanks{Corresponding author} }  \date{\today}  \begin{document}  \maketitle \thispagestyle{empty}  \begin{abstract} We investigate whether persistent latent state modulation can achieve functional causal control in frozen pretrained LLMs. Using PLAA (Persistent Latent Adaptive Architecture)---a multiplicative FiLM mechanism with GRU-driven state evolution applied to frozen Qwen3-4B---we find a surprising dissociation: latent modulation induces genuine structural differentiation across layers ($6.6\times$ growth in inter-layer weight variance) but contributes only $\Delta \approx 0.04$ to inference-time language modeling, with persistent state coupling of exactly zero.  However, the same mechanism provides substantial \textbf{training-time regularization}: LoRA-only training on multi-turn trajectory data diverges (loss $= 7.81$, worse than the frozen baseline of $4.40$), while the full modulation system converges to $3.41$. The LoRA weights trained within the modulation framework retain their quality even when FiLM is removed at inference.  These results establish that \textbf{inference-time functional coupling} and \textbf{training-time optimization shaping} are separable phenomena in modulated LLM residual streams. Latent modulation can reshape the optimization landscape ($\partial \mathcal{L} / \partial W$ is substantially altered) without participating in the inference computation ($\partial y / \partial S_t \approx 0$). We additionally document a critical methodological failure: PEFT's \texttt{save\_pretrained()} silently discards custom injection layers, invalidating ablation studies that do not independently verify checkpoint weights. \end{abstract}  \newpage \setcounter{page}{1}  %------------------------------------------------------------------ \section{Introduction} %------------------------------------------------------------------  Recent work on recurrent memory, latent-variable dialogue, and activation engineering shares a common premise: injecting persistent signals into transformer residual streams can modulate model behavior in a controllable, state-dependent manner. This premise is rarely tested in isolation from the training objective that produces it.  \textbf{The question.} Does persistent latent modulation in frozen pretrained LLMs achieve functional causal control at inference time? Or does the modulation signal serve a different role?  \textbf{This paper.} We present a systematic empirical study using PLAA (Persistent Latent Adaptive Architecture), testing three injection mechanisms---additive gated cross-attention, AdaLN FiLM, and multiplicative FiLM---in a frozen Qwen3-4B model. Our central finding is a three-way dissociation:  \begin{enumerate}[leftmargin=*, label=\textbf{\arabic*.}]  \item \textbf{Structural symmetry breaking is real.} Multiplicative FiLM induces $6.6\times$ growth in inter-layer modulation variance, with rich phase dynamics including regime-hopping and metastable hierarchy formation.  \item \textbf{Functional causal control is absent.} Controlled ablation shows inference-time $\Delta \approx 0.04$ (reproduced across two independent training runs). Freezing the latent state produces $\Delta = 0.000$.  \item \textbf{Training-time scaffold effect is strong.} The modulation mechanism combined with auxiliary losses prevents LoRA divergence on trajectory data, enabling convergence where LoRA-only training fails catastrophically.  \end{enumerate}  The mechanistic interpretation is that latent modulation reshapes the optimization landscape without participating in inference computation: $\partial \mathcal{L} / \partial W$ is substantially altered by the modulation structure, but $\partial y / \partial S_t \approx 0$.  \textbf{Additional contributions:} (a) a taxonomy of injection failure modes (additive absorption, AdaLN destructive coupling); (b) documentation of a critical PEFT checkpoint failure that silently invalidates ablation studies; (c) a corrected evaluation protocol for any research combining PEFT adapters with custom injection layers.  %------------------------------------------------------------------ \section{Related Work} %------------------------------------------------------------------  \textbf{Latent-variable dialogue.} VHRED and related models maintain latent variables across dialogue turns but operate on trainable (non-frozen) decoders. PLAA tests whether similar mechanisms can modulate frozen pretrained LLMs.  \textbf{Parameter-efficient modulation.} LoRA, IA3, prefix tuning, and adapter layers all inject learned parameters into frozen transformers. These methods optimize for task performance directly. PLAA differs in maintaining a \emph{persistent dynamical state} that evolves across turns, testing whether temporal dynamics translate to functional control.  \textbf{Activation engineering and steering vectors.} Recent work demonstrates that additive interventions on residual streams can steer model behavior. Our additive injection results (Section~\ref{sec:additive}) are consistent with this literature when interventions are large, but show that small learned additive signals are absorbed under the tested regime.  \textbf{FiLM and conditional modulation.} Feature-wise Linear Modulation originates in visual question answering. We adapt it to LLM residual streams and find that its primary contribution is to training dynamics rather than inference-time conditioning.  %------------------------------------------------------------------ \section{Method} %------------------------------------------------------------------  \subsection{Architecture}  PLAA injects a learnable latent state $\St \in \R^{256}$ into a frozen Qwen3-4B model (4-bit NF4 quantization). The state evolves via a GRU cell: \[ \Stp = \mathrm{GRU}(\St,\, h_{\mathrm{last}})$
where $h_{\mathrm{last}}$ is the final non-padded hidden state after the forward
pass. Modulation is applied to layers 16--28 (13 layers total).

\subsubsection*{Multiplicative FiLM (Primary Mechanism)}

\begin{lstlisting}
Each injected layer l in [16, 28]:
  h_res  = self_attention(h)           # frozen
  scale  = tanh(W_l . S_t)            # per-layer, W_l in R^{2560x256}
  h_out  = h_res * (1 + alpha * scale) # alpha = 0.1
  h_out  = MLP(h_out)                 # frozen
\end{lstlisting}

Zero initialization ensures modulation starts as identity. Trainable parameters:
17.5M (13 scale projections + GRU + LoRA adapters).

\subsection{Training}

\subsubsection*{Data}

3,197 dialogue micro-trajectories (3--5 turns each): 60\% PersonaChat, 30\% ESConv,
10\% Bluemoon.

\subsubsection*{Loss}

$\mathcal{L} = \underbrace{0.3\cdot\mathcal{L}_\mathrm{CE}}_\text{language model} + \underbrace{0.5\cdot\|\Delta S\| + 0.3\cdot\mathcal{L}_\mathrm{repulsion} + 0.3\cdot\|S\|}_\text{auxiliary regularization} + \underbrace{(-0.001)\cdot\mathrm{Var}(\|W_\ell\|)}_\text{diversity pressure}$

The auxiliary terms (coherence, repulsion, norm) regularize state dynamics. The
diversity term encourages inter-layer specialization.

\subsubsection*{Hyperparameters}

Optimizer: AdamW, lr $= 10^{-4}$, batch size 2. Training: 2,650--2,850 steps
depending on run. GPU: NVIDIA RTX PRO 6000 Blackwell (96\,GB VRAM).

\subsection{Evaluation Protocol}
\label{sec:protocol}

\textbf{Checkpoint correction.} PEFT's \texttt{save\_pretrained()} persists only
LoRA weights, silently discarding custom injection layers. This failure is
undetectable without explicit weight verification: zero-initialized FiLM produces
identity modulation, so the model functions normally with incorrect weights.

All results in this paper use independently saved and verified FiLM weights:

\begin{lstlisting}
# Save
torch.save({"plaa_core": plaa_core.state_dict(),
            "scale_proj": {i: blocks[i].scale_proj.state_dict()
                           for i in range(16, 29)}},
           f"{ckpt}/plaa_core.pt")

# Verify after loading
scale_norms = sum(l.scale_proj.weight.norm() for l in layers)
assert scale_norms > 1.0, "Weights are zero-initialized!"
\end{lstlisting}

\textbf{Ablation conditions.} All causal measurements use pure LM cross-entropy
(no auxiliary terms) on held-out prompts:
\begin{itemize}
  \item \textbf{Trained:} full system with normal state evolution
  \item \textbf{FiLM removed:} scale forced to zero (modulation disabled)
  \item \textbf{State frozen:} $S_t = S_0$ for all turns (GRU disabled)
\end{itemize}

%------------------------------------------------------------------
\section{Results}
%------------------------------------------------------------------

\subsection{Central Result: Training-Time Scaffold}
\label{sec:scaffold}

\begin{table}[h]
\centering
\begin{tabular}{lccc}
\toprule
\textbf{Condition} & \textbf{Pure LM Loss} & \textbf{vs.\ Vanilla} & \textbf{Converged?} \\
\midrule
Vanilla Qwen3-4B (frozen)         & 4.403 & ---     & --- \\
LoRA-only (same data, same steps) & 7.810 & $+3.41$ & \textbf{No} \\
mFiLM trained (full system)       & 3.413 & $-0.99$ & Yes \\
mFiLM, FiLM removed at inference  & 3.379 & $-1.02$ & Yes \\
\bottomrule
\end{tabular}
\caption{Training-time scaffold effect. LoRA-only diverges on multi-turn trajectory
data despite using all-token cross-entropy loss. The mFiLM system converges, and
the trained LoRA weights retain quality when FiLM is removed at inference.}
\end{table}

\textbf{Interpretation.} LoRA adapters trained within the FiLM + auxiliary loss
framework reach a qualitatively different optimum than LoRA trained alone. The
FiLM mechanism does not contribute at inference (next section), but its presence
during training prevents divergence. The auxiliary losses (coherence, norm,
repulsion) provide gradient signals that stabilize LoRA optimization on multi-turn
trajectory data where pure cross-entropy is insufficient.

\subsection{Causal Ablation: Inference-Time Contribution}

\begin{table}[h]
\centering
\begin{tabular}{lccc}
\toprule
\textbf{Condition} & \textbf{v1 Loss} & \textbf{v2 Loss} & $\Delta$ \\
\midrule
Trained (normal $S$)     & 2.700 & 3.413 & --- \\
FiLM removed (scale $=0$) & 2.744 & 3.379 & $+0.044$ / $+0.035$ \\
State frozen ($S = S_0$) & 2.700 & 3.413 & $+0.000$ / $+0.000$ \\
\bottomrule
\end{tabular}
\caption{Causal ablation across two independent training runs. v1: single-token
training loss (2,650 steps). v2: all-token CE (2,850 steps). FiLM causal
$\Delta \approx 0.04$ is stable. State coupling is zero in both.}
\end{table}

The causal contribution of FiLM at inference is $\Delta \approx 0.04$---real but
negligible relative to the $0.99$--$1.02$ improvement over vanilla. The persistent
state $\St$ is computationally inert: freezing it produces no measurable change.

\subsection{Structural Dynamics}

Despite functional inertness, the modulation system exhibits rich structural
dynamics during training:

\begin{table}[h]
\centering
\begin{tabular}{ccc}
\toprule
\textbf{Step} & $w_\mathrm{std}$ & \textbf{cos (inter-layer)} \\
\midrule
   50 & 0.062 & 0.61 \\
  350 & 0.244 & 0.35 \\
  800 & 0.333 & 0.52 \\
 1800 & 0.356 & 0.63 \\
 2650 & 0.404 & 0.48 \\
\bottomrule
\end{tabular}
\caption{Layer specialization growth (v1). $6.6\times$ increase in inter-layer
modulation variance over training. See Appendix~\ref{app:dynamics} for full
phase analysis.}
\end{table}

Layers develop distinct modulation profiles ($w_\mathrm{std}: 0.06 \to 0.40$),
but this structural differentiation does not translate to functional specialization
at inference time.

\subsection{Preliminary Architectures: Additive and AdaLN}
\label{sec:additive}

Before arriving at multiplicative FiLM, two alternative injection mechanisms were
tested and found insufficient:

\begin{table}[h]
\centering
\begin{tabular}{lccp{5.5cm}}
\toprule
\textbf{Mechanism} & \textbf{Causal} $\Delta$ & \textbf{LM Loss} & \textbf{Failure Mode} \\
\midrule
Additive (gated CA) & 2.01 (vs.\ noise $\pm 19$) & --- & Absorbed by residual stream \\
AdaLN FiLM          & $+10.50$ (harmful) & 13.38 & Destructively coupled to aux.\ loss \\
Mult.\ FiLM         & $+0.04$ & 3.41 & Structurally active, functionally inert \\
\bottomrule
\end{tabular}
\caption{Architecture comparison. Additive injection is below noise floor. AdaLN
is causally active but harmful. Multiplicative FiLM is the only viable architecture.
See Appendix~\ref{app:additive} for detailed ablations.}
\end{table}

%------------------------------------------------------------------
\section{Analysis}
%------------------------------------------------------------------

\subsection{Inference Graph vs.\ Optimization Graph}

The central finding admits a precise mechanistic characterization. Consider two
computational graphs:

\textbf{Inference graph:} $x \to h_{\mathrm{layers}} \to y$, where modulation
applies $h \cdot (1 + 0.1 \cdot \tanh(W_\ell \cdot S_t))$. With $\alpha = 0.1$
and $\tanh$ bounding:
$\left\|\frac{\partial y}{\partial S_t}\right\| \leq 0.1 \cdot \|h\| \cdot \|W_\ell\| \cdot (1 - \tanh^2(W_\ell S_t))$
This gradient is bounded and, empirically, negligible: the frozen residual stream
compensates for $\pm 10\%$ perturbations within its natural operating range.

\textbf{Optimization graph:} $\theta_{\mathrm{LoRA}}, W_\ell, S_0 \to
\mathcal{L}_\mathrm{total}$. Here the modulation structure participates through:
\begin{itemize}
  \item Auxiliary loss gradients ($\nabla_{W_\ell} \mathcal{L}_\mathrm{aux}$) that
    provide non-zero gradient signal even when LM gradients are noisy
  \item Multiplicative coupling between $W_\ell$ and LoRA parameters that creates
    implicit gradient regularization
  \item GRU state evolution that enforces temporal smoothness constraints on the
    optimization trajectory
\end{itemize}

The result: $\partial \mathcal{L} / \partial W_{\mathrm{LoRA}}$ is substantially
different in the presence of FiLM + auxiliary losses, guiding LoRA to a better
minimum. But $\partial y / \partial S_t \approx 0$: the modulation does not
participate in the final computation.

\subsection{Why LoRA-Only Diverges}

Multi-turn trajectory data presents a challenging optimization landscape for
LoRA alone:
\begin{itemize}
  \item Each training sample spans 3--5 dialogue turns with state dependencies
  \item Pure cross-entropy provides sparse, high-variance gradients across turns
  \item Without auxiliary regularization, LoRA adapters overfit to local token
    patterns while losing global coherence
\end{itemize}

The FiLM system addresses this through multiple mechanisms: the coherence loss
($\|\Delta S\|$) penalizes abrupt representational shifts between turns; the norm
penalty ($\|S\|$) prevents state explosion; the diversity pressure encourages
layer-wise specialization that distributes the representational burden. These
signals flow through the shared computation graph to regularize LoRA gradients.

\subsection{Bounded Modulation and Residual Stream Robustness}

The $\alpha = 0.1$ operating point creates a fundamental ceiling on inference-time
influence. With hidden states $\|h\| \sim 50$--$100$, the modulation perturbation
is bounded by $\pm 10\%$ of the residual stream magnitude. The frozen transformer
layers downstream compensate for perturbations within this range, effectively
absorbing the modulation signal.

Alpha ablation (Appendix~\ref{app:alpha}) confirms that increasing $\alpha$ beyond
$0.1$ produces harmful interference rather than stronger control: the LoRA adapters
were optimized under $\alpha = 0.1$ and cannot adapt to larger modulation without
retraining.

\subsection{Structure--Function Dissociation}

Weight-space metrics ($w_\mathrm{std}$, gradient cosine similarity) can be
misleading indicators of functional influence. The multiplicative FiLM system
exhibits:
\begin{itemize}
  \item $6.6\times$ growth in inter-layer weight variance
  \item Phase transitions and regime-hopping dynamics
  \item Metastable coupled hierarchy formation
\end{itemize}

None of these structural phenomena translate to measurable functional
differentiation at inference time. This dissociation implies that structural
complexity in modulation weights is necessary for the training-time scaffold
effect (the diversity pressure requires weight differentiation) but is not
sufficient for inference-time causal control.

%------------------------------------------------------------------
\section{Limitations}
%------------------------------------------------------------------

\begin{enumerate}
  \item \textbf{Single base model.} All experiments use Qwen3-4B under NF4
    quantization. The robustness of frozen residual streams to modulation may
    differ across model families and scales.

  \item \textbf{No comparison to existing modulation methods.} We do not compare
    against prefix tuning, IA3, soft prompts, or activation steering under matched
    conditions. Future work should establish whether these methods exhibit similar
    structure--function dissociation.

  \item \textbf{Fixed $\alpha$ regime.} The $\alpha = 0.1$ bound was chosen for
    training stability. A purpose-designed curriculum that gradually increases
    $\alpha$ may achieve stronger coupling.

  \item \textbf{Evaluation scope.} Causal ablation measures pure LM cross-entropy.
    Task-specific evaluations (emotion recognition, persona consistency) may reveal
    functional contributions not captured by perplexity.

  \item \textbf{Training data scale.} 3,197 trajectories is small by modern
    standards. Larger-scale training may produce different dynamics.
\end{enumerate}

%------------------------------------------------------------------
\section{Conclusion}
%------------------------------------------------------------------

We set out to build persistent state-driven modulation for frozen LLMs and found
something different: a training-time optimization scaffold.

Multiplicative FiLM modulation in frozen Qwen3-4B achieves genuine structural
symmetry breaking ($6.6\times$ inter-layer variance growth) but contributes only
$\Delta \approx 0.04$ at inference, with zero state coupling. The persistent latent
state---the core design premise---is computationally inert. Yet the same mechanism,
combined with auxiliary losses, prevents LoRA divergence on trajectory data and
enables convergence to below-vanilla performance.

The mechanistic explanation is a dissociation between the inference graph and the
optimization graph: latent modulation reshapes $\partial \mathcal{L} / \partial
W_{\mathrm{LoRA}}$ without participating in $\partial y / \partial S_t$. This
establishes that structural differentiation, functional causal control, and
training-time regularization are three separable phenomena in modulated LLM
residual streams.

\textbf{Implications.} Any work claiming functional latent control must verify:
(1) checkpoint weights are non-zero; (2) causal ablation uses pure task loss;
(3) state coupling is tested via freeze ablation. Structural metrics alone are
insufficient evidence of functional influence.

\textbf{Future directions.} Three approaches may bridge the structure--function gap:
unbounded modulation with curriculum training; gated routing that makes the state
causally necessary for prediction; and mixture-of-experts architectures where
$\St$ selects computation paths rather than perturbing residual magnitudes.

%------------------------------------------------------------------
\appendix
%------------------------------------------------------------------

\section{Additive Injection: Detailed Ablation}
\label{app:additive}

\subsection*{Architecture}

Gated cross-attention: $h_s + \mathrm{gate} \cdot e_o$, where
$e_{\mathrm{mem}} = \mathrm{proj}(\St) \in \R^{16 \times 2560}$,
gate initialized $0.1$ with $\tanh$ activation. 387M trainable parameters.

\subsection*{Multi-Stage Attenuation}

Three compounding factors reduce effective influence:
\begin{enumerate}
  \item $\|S_t\| \approx 0.15$ (norm penalty) $\Rightarrow$ small projection
  \item $\tanh(0.1) \approx 0.1$ gate scale
  \item $\times 13$ residual layers $\Rightarrow$ effective $\approx 0.0003 \cdot \|h\|$
\end{enumerate}

\subsection*{Gate Ablation}

\begin{table}[h]
\centering
\begin{tabular}{lcc}
\toprule
\textbf{Condition} & \textbf{PPL} & $\Delta$\textbf{PPL} \\
\midrule
gate $= 0$      & 27.03 & ---  \\
gate $= 1.0$    & 26.26 & 0.77 \\
gate $= 2.0$    & 25.02 & 2.01 \\
gate $=$ random & 25.94 & 1.09 \\
\bottomrule
\end{tabular}
\caption{Max $\Delta = 2.01$, noise std $= \pm 19$. Below noise floor.}
\end{table}

\section{AdaLN FiLM: Detailed Analysis}
\label{app:adaln}

\subsection*{Architecture}

$\mathrm{LN}(h_s) \cdot (1 + \alpha\cdot\tanh(W_\mathrm{scale}\,\St)) +
\alpha\cdot\tanh(W_\mathrm{shift}\,\St)$, with $\alpha = 0.1$. 17.5M parameters.

\subsection*{Invalid vs.\ Corrected Measurement}

\begin{table}[h]
\centering
\begin{tabular}{lccl}
\toprule
& \textbf{Invalid (composite)} & \textbf{Corrected (pure LM)} & \textbf{Interpretation} \\
\midrule
Trained       & 14.68 & 13.382 & System active \\
FiLM removed  & 3.14  & 2.878  & Better without FiLM \\
$\Delta$      & 11.44 & $+10.50$ & \textbf{Harmful} \\
State frozen  & 13.35 & 13.382 & Zero coupling \\
\bottomrule
\end{tabular}
\caption{AdaLN modulation is destructively coupled to the training objective.
The trained system (13.38) is far worse than vanilla Qwen3-4B (3.70).}
\end{table}

\subsection*{Temporal--Specialization Tension}

\begin{table}[h]
\centering
\begin{tabular}{lccc}
\toprule
\textbf{Variant} & $w_\mathrm{std}$ & $|S|$ & \textbf{cos} \\
\midrule
No coupling (v5)       & 0.033 & 0.57 & 0.15 \\
Coupling 0.01 (v6)     & 0.000 & 2.59 & 0.19 \\
Coupling 0.001 (v7)    & 0.000 & 0.44 & 0.16 \\
\bottomrule
\end{tabular}
\caption{Even weak coupling ($0.001\times$) collapses layer specialization.}
\end{table}

\section{Phase Dynamics and Regime-Hopping}
\label{app:dynamics}

\subsection*{Four-Phase Taxonomy}

\begin{table}[h]
\centering
\begin{tabular}{clccc}
\toprule
\textbf{Phase} & \textbf{Name} & $w_\mathrm{std}$ & \textbf{cos} & \textbf{Dynamics} \\
\midrule
I   & Symmetric          & $<0.05$       & $>0.9$      & Homogeneous \\
II  & Convergent         & $0.05$--$0.3$ & $0.3$--$0.7$& Emerging specialization \\
III & Specialized        & $0.3$--$0.5$  & $0.3$--$0.5$& Stable hierarchy \\
IV  & Metastable Coupled & $0.35$--$0.5$ & $0.4$--$0.7$& Hierarchy + bursts \\
\bottomrule
\end{tabular}
\caption{Phase taxonomy based on structural metrics. Multiplicative FiLM (v1)
reaches Phase IV. Phase classification does not imply functional specialization.}
\end{table}

\subsection*{Regime-Hopping Events (v1)}

Two synchronization bursts (step 1750: cos $= 0.627$; step 1950: cos $= 0.696$)
briefly restore global coupling. $w_\mathrm{std}$ does not collapse during bursts
($0.348 \to 0.382$), and loss recovers to lower baseline afterward. Structural
differentiation strengthens under perturbation.

\subsection*{v2 State Dynamics (All-Token CE)}

Under all-token CE training, state dynamics are qualitatively different from v1:

\begin{table}[h]
\centering
\begin{tabular}{cccc}
\toprule
& \textbf{v1 (single-token)} & \textbf{v2 (all-token)} & \textbf{Ratio} \\
\midrule
$|S|$ baseline    & 0.15--0.57 & 2.2--3.9  & $\sim 7\times$ \\
$|S|$ burst       & not observed & 9.13     & --- \\
$dS$ range        & 0.5--1.0   & 0.5--8.6  & $\sim 5\times$ \\
$w_\mathrm{std}$ @ step 1200 & $\sim 0.30$ & 0.09--0.10 & $0.3\times$ \\
\bottomrule
\end{tabular}
\caption{v2 shows more active state dynamics but slower layer specialization,
suggesting all-token CE prioritizes state evolution over structural differentiation.
Despite these differences, final causal $\Delta$ is equivalent ($\approx 0.04$).}
\end{table}

The v2 system exhibits regime-hopping bursts ($|S| = 9.13$ at step 1050, recovering
by step 1150) not observed in v1. This confirms that all-token CE provides stronger
gradient signal to the GRU, producing more active state dynamics. However, this
increased dynamism does not translate to increased functional coupling---the final
state freeze $\Delta$ remains zero.

\section{Alpha Ablation and Causal Sensitivity}
\label{app:alpha}

Starting from the multiplicative FiLM checkpoint at step 850 ($w_\mathrm{std} =
0.307$), we swept $\alpha \in \{0.1, 0.5, 1.0\}$:

\begin{table}[h]
\centering
\begin{tabular}{lcccc}
\toprule
$\alpha$ & \textbf{Step} & $\Delta$ \textbf{causal} & $w_\mathrm{std}$ & \textbf{Loss} \\
\midrule
\multirow{2}{*}{0.1} & 0   & $-0.045$ & 0.307model to \emph{depend} on
    the modulation signal for correct next-token prediction, making the latent state
    causally necessary rather than optional.
  \item \textbf{Competitive computation:} Mixture-of-experts style routing where
    $\St$ selects which expert processes each token, creating a hard dependency on
    the state.
\end{itemize}

The central open problem is clear: how to convert structural differentiation into
functional specialization, and how to make the latent state causally necessary rather
than decorative.

%------------------------------------------------------------------
\appendix
%------------------------------------------------------------------

\section{Full Training Pipeline}

\begin{table}[h]
\centering
\begin{tabular}{llccc}
\toprule
\textbf{Phase} & \textbf{Content} & \textbf{Steps} & \textbf{Params} & \textbf{Status} \\
\midrule
1              & YaRN 256K context extension & ---  & ---    & Complete \\
2              & SFT (134K samples)          & ---  & ---    & Complete \\
2.5            & Persona alignment           & ---  & ---    & Complete \\
3A             & Stability validation        & 1K   & ---    & Complete \\
3B             & Gradient flow verification  & ---  & ---    & Complete \\
3C v1          & Single-step additive        & 5K   & 387M   & Complete \\
3C v2          & Trajectory additive         & 1.4K & 387M   & Complete \\
3C v3          & Additive + gate noise       & 6.7K & 387M   & Baseline \\
3C FiLM (AdaLN)& AdaLN modulation           & 7K   & 17.5M  & Ablated  \\
\textbf{3C mFiLM} & \textbf{Multiplicative FiLM} & \textbf{2.65K} & \textbf{17.5M} & \textbf{Primary} \\
\bottomrule
\end{tabular}
\caption{Full training pipeline.}
\end{table}

\section{Scale/Shift Statistics (AdaLN)}

\begin{table}[h]
\centering
\begin{tabular}{lcc}
\toprule
\textbf{Metric} & \textbf{Step 1000} & \textbf{Step 2000} \\
\midrule
scale mean        & 0.0001 & 0.0001 \\
scale std         & 0.0060 & 0.0055 \\
scale max         & 0.0228 & 0.0209 \\
layer\_scale\_std & 0.0060 & 0.0055 \\
\bottomrule
\end{tabular}
\caption{AdaLN scale/shift statistics. All 13 layers produce nearly identical
modulation, confirming no layer specialization under baseline AdaLN training.}
\end{table}

\section{Alpha Ablation and Causal Sensitivity}
\label{app:alpha}

Starting from the multiplicative FiLM checkpoint at step 850 ($w_\mathrm{std} =
0.307$), we swept $\alpha \in \{0.1, 0.5, 1.0\}$, measuring causal $\Delta$ at
load time (step 0) and after 500 additional training steps.

\begin{table}[h]
\centering
\begin{tabular}{lcccc}
\toprule
$\alpha$ & \textbf{Step} & $\Delta$ \textbf{causal} & $w_\mathrm{std}$ & \textbf{Loss (trained)} \\
\midrule
\multirow{2}{*}{0.1} & 0   & $-0.045$ & 0.307 & 2.700 \\
                     & 500 & $+0.095$ & 0.515 & 2.682 \\
\midrule
\multirow{2}{*}{0.5} & 0   & $+0.238$ & 0.307 & 2.938 \\
                     & 500 & $+0.526$ & 0.539 & 3.225 \\
\midrule
\multirow{2}{*}{1.0} & 0   & $+2.679$ & 0.307 & 5.381 \\
                     & 500 & $+2.862$ & 0.467 & 6.621 \\
\bottomrule
\end{tabular}
\caption{Alpha ablation. Positive $\Delta$ indicates removing FiLM \emph{improves}
loss (modulation is harmful). All values are pure LM cross-entropy.}
\end{table}

Three findings emerge:

\textbf{1.} Causal sensitivity grows with $\alpha$: $-0.045 \to +2.862$.

\textbf{2.} Strong modulation ($\alpha \geq 0.5$) is harmful. The frozen LoRA adapters and Qwen backbone were optimized under $\alpha = 0.1$; abruptly increasing $\alpha$ disrupts learned compensation patterns, and 500 additional training steps fail to recover.

\textbf{3.} The $\alpha = 0.1$ operating point is real but narrow. The paper's central finding---structural differentiation without proportional functional influence---holds within this regime. The narrow operating window represents a genuine architectural limitation.

\begin{thebibliography}{9}
\bibitem{perez2018film} Perez, E., et al. FiLM: Visual Reasoning with a General Conditioning Layer. \textit{AAAI 2018}.
\bibitem{peebles2023scalable} Peebles, W. \& Xie, S. Scalable Diffusion Models with Transformers. \textit{ICCV 2023}.
\bibitem{huang2017arbitrary} Huang, X. \& Belongie, S. Arbitrary Style Transfer in Real-time with Adaptive Instance Normalization. \textit{ICCV 2017}.
\bibitem{sukhbaatar2019augmenting} Sukhbaatar, S., et al. Augmenting Self-attention with Persistent Memory. \textit{arXiv:1907.01470}.
\bibitem{wu2022memorizing} Wu, Y., et al. Memorizing Transformers. \textit{ICLR 2022}.
\end{thebibliography}

\end{document}
